\section{Conclusion and Future Work}\label{sec:conclusion}
We presented PHQD, a modular patch-based hybrid quantum--classical pipeline that removes Gaussian noise from the latent maps of a convolutional autoencoder on Fashion-MNIST. Overlapping $2\times2$ patches are denoised by hybrid models that combine classical projections with a parameterized quantum circuit, and the patches are averaged back into a latent map. We compared four 4-qubit configurations and a 6-qubit model with data re-uploading, three-basis readout and joint stitching refinement against classical MLP and U-Net baselines.

All denoisers improved on the noisy input, but the quantum models did not beat the classical baselines. On 32 test images, the MLP gave the best image MSE (0.00245, against 0.00294 for the best 4-qubit model, EPQD-L, and 0.00416 for the noisy input), and the U-Net gave the best FID (51.2, against about 70 for the EPQD models). Circuit depth and the IQP-style encoding changed the results only slightly. These results show why strong classical baselines are needed before crediting an improvement to the circuit.

Future work includes larger test sets, repeated runs with several random seeds, other datasets and noise levels, a parameter-matched classical replacement for the circuit, and a complete ablation of the 6-qubit model. We also plan hardware-aware simulation with realistic noise and, later, tests on real quantum devices.
